%!TEX root = ../../../standalone/applications/app-scimba-hybrid.tex
\ApplicationFicheHeading{app-scimba-hybrid}
\textbf{app-scimba-hybrid: Hybrid Numerical Methods}
\label{sec:app:specs:app-scimba-hybrid}

% Generated by scripts/sync_applications.py; edit the CSV snapshot.
\begin{FicheMetadataTable}
id & app-scimba-hybrid \\
name & Hybrid  numerics \\
Partners & Unistra \\
PC & PC1 - ExaMA \\
Responsible (Permanent) & Emmanuel Franck \\
WP7 Engineer & \todo[app-scimba-hybrid / WP7 Engineer]{Sheet field missing.} \\
work\_package & WP2 \\
application\_type & demonstrator \\
purpose & Test numerical hybrid methods on high dimensional physical and parametric PDE problems \\
inputs & \todo[app-scimba-hybrid / inputs]{Sheet field missing.} \\
outputs & \todo[app-scimba-hybrid / outputs]{Sheet field missing.} \\
metrics & \todo[app-scimba-hybrid / metrics]{Sheet field missing.} \\
status & \todo[app-scimba-hybrid / status]{Sheet field missing.} \\
Benchmark scope & \todo[app-scimba-hybrid / Benchmark scope]{Sheet field missing.} \\
Framework & Scimba \\
parallel\_framework & GPU \\
repo\_url & \todo[app-scimba-hybrid / repo\_url]{Sheet field missing.} \\
tex\_url & \todo[app-scimba-hybrid / tex\_url]{Sheet field missing.} \\
\end{FicheMetadataTable}


\subsection{Description of the benchmark}
\textbf{Status:} Proposed | \textbf{Partners:} Unistra | \textbf{Responsible:} Emmanuel Franck

\textbf{Purpose:} Test numerical hybrid methods on high-dimensional physical and parametric PDE problems (demonstrator).

\textbf{Key Methods:}
\begin{itemize}
\item WP2: Hybrid numerical-ML methods, high-dimensional PDEs
\end{itemize}

\textbf{Framework:} Scimba

\textbf{Scientific Context:} Hybrid approaches combine traditional numerical methods (FEM, FV, spectral) with machine learning components. This can provide the best of both worlds: guaranteed conservation properties and numerical stability from classical methods, with efficiency and adaptivity from ML.

\subsection{Benchmarking methodology}
\begin{description}
  \item[Framework:] Scimba
  \item[parallel\_framework:] GPU
\end{description}
\todo[app-scimba-hybrid / benchmarking methodology]{Benchmarking methodology to be supplied by contributors.}

\subsection{Input/Output Dataset Description}

\subsubsection{Input Data}
\todo[app-scimba-hybrid / inputs]{Sheet field missing.}

\subsubsection{Output Data}
\todo[app-scimba-hybrid / outputs]{Sheet field missing.}

\subsection{Results summary}
\todo[app-scimba-hybrid / results summary]{Results summary to be supplied by contributors.}
